% ECONOMICS_PROBLEM: {"id": "C73-25", "title": "Minimizer selection by the uncorrected vanishing-common-noise game", "jel_code": "C73", "source_id": "OP-0215"}
% FIRST_POSTED: 2026-10-09
% ATTEMPT_STATUS: unresolved
% ENTRY_KIND: research_attempt
% !TeX program = pdflatex
% Self-contained source: pdflatex twice; no fontspec or external bibliography.
\documentclass[11pt,a4paper]{article}
\usepackage[margin=24mm,headheight=14pt]{geometry}
\usepackage[T1]{fontenc}
\usepackage[utf8]{inputenc}
\usepackage{lmodern,microtype}
\usepackage{amsmath,amssymb,mathtools,amsthm}
\usepackage{enumitem,booktabs,longtable,array,xurl,listings,needspace}
\usepackage[hidelinks,unicode]{hyperref}
\hypersetup{pdftitle={Economic Theory: Proof Attempts and Adversarial Checks},pdfauthor={Research notes},pdfsubject={Eleven user-supplied mathematical problems}}
\newcommand{\NE}{\operatorname{NE}}
\DeclareMathOperator*{\argmax}{arg\,max}
\DeclareMathOperator*{\argmin}{arg\,min}
\newtheorem{theorem}{Theorem}[section]
\newtheorem{lemma}[theorem]{Lemma}
\newtheorem{proposition}[theorem]{Proposition}
\newtheorem{corollary}[theorem]{Corollary}
\theoremstyle{definition}
\newtheorem{definition}[theorem]{Definition}
\theoremstyle{remark}
\newtheorem*{remark}{Remark}
\setlength{\parindent}{0pt}
\setlength{\parskip}{5pt plus 1pt minus 1pt}
\setlength{\emergencystretch}{2em}
\setcounter{tocdepth}{1}
\setlist{leftmargin=*,itemsep=2pt,topsep=4pt}
\lstset{basicstyle=\small\ttfamily,columns=fullflexible,keepspaces=true,
 breaklines=true,breakatwhitespace=false,showstringspaces=false,
 frame=single,framerule=0.25pt,aboveskip=8pt,belowskip=8pt,xleftmargin=3pt,xrightmargin=3pt}
\allowdisplaybreaks[1]
\raggedbottom

\begin{document}
{\large\bfseries Research attempt and adversarial verification}\par
9 October 2026\hfill Original-target status: \textbf{UNRESOLVED}\par
\section{C73-25: Uncorrected vanishing-common-noise potential selection}
\label{sec:C73-25}
\noindent\textbf{Input:} \texttt{C73-25.tex}. \quad\textbf{Tools:} web browsing [YES]; Python computation [YES].
\subsection*{1) FORMAL RESTATEMENT}
Let $d\geq3$, $T>0$, $\Delta_d=\{m\in[0,1]^d:\sum_i m_i=1\}$, and $m_0\in\operatorname{int}\Delta_d$. Smooth potentials $\mathcal F,\mathcal G$ are defined on a neighborhood of the simplex; set $f_i=\partial_i\mathcal F$, $g_i=\partial_i\mathcal G$. Write
\[
 H_i(U)=\tfrac12\sum_{j\ne i}(U_i-U_j)_+^2,
 \quad b_i(m,U)=\sum_{j\ne i}[m_j(U_j-U_i)_+-m_i(U_i-U_j)_+],
\]
\[
 R_i^\phi(m)=\sum_{j\ne i}[m_j\phi_\varepsilon(m_i)-m_i\phi_\varepsilon(m_j)],
 \quad L_\varepsilon=\varepsilon\sum_{j,k}(m_j\delta_{jk}-m_jm_k)\partial_{jk}^2.
\]
A family $\phi_\varepsilon$ is smooth, nonnegative and locally uniformly vanishing on $(0,1]$. It is admissible in this note precisely when the input's following PDE has a unique solution in its specified classical class and the associated SDE has a unique interior-valued solution:
\begin{align*}
0={}&\partial_tU_i+L_\varepsilon U_i+(b+R^\phi)\cdot DU_i
 +2\varepsilon(e_i-m)\cdot DU_i-H_i(U)+f_i(m)\\
&+\sum_{j\ne i}\phi_\varepsilon(m_j)(U_j-U_i),
 \qquad U_i(T,m)=g_i(m),\\
dm_i^\varepsilon={}&[b_i(m^\varepsilon,U^\varepsilon)+R_i^\phi(m^\varepsilon)]dt
 +\sqrt\varepsilon\sum_j\sqrt{m_i^\varepsilon m_j^\varepsilon}
       (dW^{ij}-dW^{ji}),\qquad m^\varepsilon(0)=m_0.
\end{align*}
The $W^{ij}$ are independent Brownian motions. All displayed differential contractions are tangent to the simplex. The classical class means continuity on the closed time--simplex domain, $C^1$ time and $C^2$ interior space regularity, bounded values and first derivatives, and locally bounded interior second derivatives, separately for each $\varepsilon$. No bound uniform in $\varepsilon$ is presumed.

For measurable nonnegative rates $a_{ij}$ with $\int_0^T\sum_{i,j\ne i}m_i a_{ij}^2dt<\infty$ and absolutely continuous simplex paths satisfying
\[
 \dot m_i=\sum_{j\ne i}(m_ja_{ji}-m_ia_{ij}),\qquad m(0)=m_0,
\]
define
\[
 J(m,a)=\int_0^T\left[\tfrac12\sum_{i,j\ne i}m_i a_{ij}^2+
 \mathcal F(m)\right]dt+\mathcal G(m(T)).
\]
Let $\mathcal M_{T,m_0}$ be the population paths occurring in global minimizers of $J$. The literal question quantifies over all the above data, every admissible repulsion family, and every sequence $\varepsilon_n\downarrow0$ with $\mathcal L(m^{\varepsilon_n})\Rightarrow Q$ in $C([0,T],\Delta_d)$: must $Q(\mathcal M_{T,m_0})=1$? The cost functions $f,g$ are \emph{not} replaced by corrected noisy costs. No correction to that target is needed. Boundary layers, the factor two in the Wright--Fisher covariance, and global versus stationary optimality are stress points.

\subsection*{2) QUICK LITERATURE/CONTEXT CHECK}
Cardaliaguet and Delarue \cite{cardaliaguet-delarue}, printed p.~3691, explicitly separate the uncorrected question from the cost-corrected selection theorem on p.~3692 (Theorem 4.5). Cecchin and Delarue \cite{cecchin-delarue} is the primary paper for that tailored potential-preserving construction. The extra costs in that theorem cannot be deleted from its hypotheses. These page numbers were checked against rendered source pages; some locators in the attachment are displaced by one page. A targeted check on 9 October 2026 did not supply a theorem covering the exact universal uncorrected quantifiers here.

\subsection*{3) ATTACK PLAN}
\textbf{Type and tools.} This is a singular-noise limit and variational selection problem. Useful tools are: direct substitution in the master equation (an exact solvable subclass); stopped SDE estimates (separate interior behavior from a boundary layer); quadratic variation and Doob's inequality (control common noise); logarithmic barriers (verify interior admissibility); vector maximum arguments (prove the zero solution is unique); kinetic-energy nonnegativity (identify global minimizers); explicit nonconvex potentials (test whether the deterministic MFG equations suffice).

\textbf{Proof tracks.} Prove a quantitative limit for zero costs with arbitrary admissible repulsion, then seek a potential-energy comparison for nonconstant costs. \textbf{Disproof tracks.} Find a nonminimizing deterministic equilibrium, test whether it is genuinely selected by the uncorrected noise, and push repulsion into a thin boundary layer. The first track succeeds for zero costs. The constructed nonminimizing equilibrium is a verified obstruction to one inference, but its noisy selection is not proved.

\subsection*{4) WORK}
\begin{theorem}[Zero-cost selection for every admissible repulsion family]
If $\mathcal F=\mathcal G=0$, then for every admissible repulsion family,
\[
 \sup_{0\leq t\leq T}\|m^\varepsilon(t)-m_0\|_\infty
 \longrightarrow0\quad\text{in probability}.
\]
Moreover $\mathcal M_{T,m_0}=\{t\mapsto m_0\}$. Consequently every subsequential path-law limit has mass one on global minimizing paths in this subclass.
\end{theorem}
\begin{proof}
The vector $U^\varepsilon\equiv0$ solves every term of the master equation and its terminal condition. Its derivatives vanish and it is in the required classical class. Admissible uniqueness therefore forces this solution, so the controlled drift $b$ is zero.

Set $\delta=\frac12\min_i m_{0,i}>0$ and
$\tau_\varepsilon=\inf\{t:\min_i m_i^\varepsilon(t)\leq\delta\}$.
Let $a_\varepsilon=\sup_{r\in[\delta,1]}\phi_\varepsilon(r)$, which tends to zero by the exact local-uniform assumption. Before $\tau_\varepsilon$, $|R_i^\phi|\leq d a_\varepsilon$; this follows by bounding separately its nonnegative incoming and outgoing sums. Write the stopped equation as
\[
 m_i^\varepsilon(t\wedge\tau_\varepsilon)-m_{0,i}
 =\int_0^{t\wedge\tau_\varepsilon}R_i^\phi(m^\varepsilon(s))ds+M_i(t).
\]
The antisymmetric Brownian representation gives
\[
 d\langle M_i\rangle_t
 =2\varepsilon m_i^\varepsilon(1-m_i^\varepsilon)
      \mathbf1_{\{t\leq\tau_\varepsilon\}}dt\leq\frac\varepsilon2dt.
\]
Doob's $L^2$ inequality ($\mathbb E\sup_{t\leq T}|M_t|^2\leq4\mathbb E|M_T|^2$ for a square-integrable martingale starting at zero) implies
$\mathbb E\sup_{t\leq T}|M_i(t)|^2\leq2\varepsilon T$.

On $\{\tau_\varepsilon\leq T\}$ some coordinate changes by at least $\delta$. For small $\varepsilon$ with $dTa_\varepsilon\leq\delta/2$, the union bound and Markov's inequality therefore give
\[
 \Pr(\tau_\varepsilon\leq T)\leq
 \Pr\left(\max_i\sup_{t\leq T}|M_i(t)|\geq\delta/2\right)
 \leq\frac{8d\varepsilon T}{\delta^2}.
\]
For any fixed $\eta>0$, also impose $dTa_\varepsilon\leq\eta/2$. On the event of no exit the full equation equals the stopped equation, and hence
\[
 \Pr\left(\sup_{t\leq T}\|m^\varepsilon(t)-m_0\|_\infty>\eta\right)
 \leq\frac{8d\varepsilon T}{\delta^2}
      +\frac{8d\varepsilon T}{\eta^2}\longrightarrow0.
\]
This proves convergence on the entire continuous-path space, not just at one time. In particular it supplies tightness and uniquely identifies every limit without assuming it in this subclass.

The deterministic objective is nonnegative kinetic energy. The zero rate matrix and constant path attain value zero. Any minimizer must satisfy $m_i a_{ij}^2=0$ almost everywhere for every edge. This implies $m_i a_{ij}=0$ almost everywhere, including where $m_i=0$. The forward equation then gives $\dot m=0$ almost everywhere. Absolute continuity and the initial condition force $m(t)=m_0$. Thus the minimizing-path set is the stated singleton.
\end{proof}

\paragraph{Boundary-layer scale in the theorem.}
No prescribed shrinking width is needed for this particular zero-cost result. The proof localizes at the \emph{fixed} interior distance $\delta=\frac12\min_i m_{0,i}$ and uses only $a_\varepsilon\to0$ there. It neither assumes nor proves that $\phi_\varepsilon$ is globally small. The probability of ever seeing the uncontrolled boundary region is explicitly $O(\varepsilon)$ after the local drift becomes small. This localization does not extend automatically when $b(m,U^\varepsilon)$ can drive the population toward that region.

\begin{proposition}[A nonvacuous admissible family for the zero-cost theorem]
For $\mathcal F=\mathcal G=0$, the family $\phi_\varepsilon(r)=2\varepsilon$ is admissible for every $T>0$ and every interior initial point. Its unique classical solution is zero. For its population law,
\[
 \mathbb E m_i^\varepsilon(t)=\frac1d+
       \left(m_{0,i}-\frac1d\right)e^{-2d\varepsilon t}.
\]
\end{proposition}
\begin{proof}
The repulsion drift is $2\varepsilon(1-dm_i)$. The SDE coefficients are locally Lipschitz in the simplex interior, preserve the sum of coordinates, and give a unique strong solution up to an exit from that interior. To exclude such an exit let $B(m)=-\sum_i\log m_i$. Using an ambient extension of this function in the tangent contractions, the generator of this population diffusion satisfies
\[
 (L_\varepsilon+R^\phi\cdot D)B
 =\varepsilon\left(\sum_i\frac1{m_i}-d\right)
  -2\varepsilon\sum_i\frac1{m_i}+2\varepsilon d^2
 \leq2\varepsilon d^2.
\]
Stop at the first coordinate below $\zeta>0$. It\^o's formula gives
$\mathbb E B(m_{T\wedge\tau_\zeta})\leq B(m_0)+2\varepsilon d^2T$.
Since $B\geq0$ on the simplex and $B\geq\log(1/\zeta)$ at an exit,
\[
 \Pr(\tau_\zeta\leq T)\leq
 \frac{B(m_0)+2\varepsilon d^2T}{\log(1/\zeta)}\longrightarrow0.
\]
Localization therefore extends the unique strong solution to all finite times and keeps it interior. Taking expectations in its bounded-coefficient SDE gives the scalar linear equation $\dot{\bar m}_i=2\varepsilon(1-d\bar m_i)$, proving the formula.

It remains to verify PDE uniqueness rather than assume it for this example. Let $V(\tau,m)=U(T-\tau,m)$ be any bounded classical solution with initial value zero, and let $M=\sup_{\tau,m,i}|V_i|$. Its forward equation is
\[
 \partial_\tau V_i=(L_\varepsilon+B_i\cdot D)V_i-H_i(V)
            +2\varepsilon\sum_{j\ne i}(V_j-V_i),
 \quad B_i=b(m,V)+2\varepsilon(\mathbf1-dm)+2\varepsilon(e_i-m).
\]
Here $\mathbf1$ is the vector with every coordinate equal to one. The vector $B_i$ is tangent. Direct calculation shows
\[
 (L_\varepsilon+B_i\cdot D)B(m)
 \leq 2M d(d-1)+\varepsilon(2d^2+d)=:C_0.
\]
Indeed, writing $a_{jk}=(V_j-V_k)_+$ gives
$b\cdot DB=\sum_{j,k\ne j}a_{jk}(1-m_j/m_k)\leq\sum_{j,k\ne j}a_{jk}\leq2Md(d-1)$.
The mutation and diffusion terms were just calculated, and the extra drift contributes $-2\varepsilon/m_i+2\varepsilon d$, which has the required upper bound.

Fix $\alpha>0$ and set $C=C_0+1$. If some component of
$V_i-\alpha[B(m)+C\tau]$ had a positive maximum on $[0,T]\times\operatorname{int}\Delta_d$, it would occur at positive time and an interior spatial point: the barrier diverges at the boundary and the initial expression is negative. At that point $V_i$ is maximal among the components. Thus $-H_i(V)\leq0$ and the switching sum is nonpositive. The spatial gradient of the subtracted expression vanishes, its diffusion contraction is nonpositive, and its time derivative at a maximum over past times is nonnegative. The PDE instead bounds that derivative by $\alpha(C_0-C)<0$, a contradiction. Therefore $V_i\leq\alpha[B+C\tau]$ everywhere. Apply the same argument to $-V_i-\alpha[B+C\tau]$: now $V_i$ is a minimum component, so $H_i(V)=0$ and the negative switching sum is nonpositive. It follows that $-V_i\leq\alpha[B+C\tau]$. Let $\alpha\downarrow0$ at each interior point, and use continuity at the boundary to obtain $V=0$ everywhere.

The zero field exists and has all the required bounded derivatives. This proves uniqueness in the stipulated class, completes admissibility, and confirms that the preceding selection theorem is not vacuous.
\end{proof}

\begin{proposition}[A smooth potential MFG equilibrium need not minimize the potential]
There is a three-state smooth potential game, with an interior initial law, having a classical deterministic equilibrium of action zero whose population path is not a global potential minimizer.
\end{proposition}
\begin{proof}
Take $T=1$, $m_0=(1/3,1/3,1/3)$, $\mathcal F=0$ and
$\mathcal G(m)=-3(m_1-1/3)^2$. Then $g_1=-6(m_1-1/3)$ and $g_2=g_3=0$. The constant path $m=m_0$, values $u=0$ and rates $a=0$ satisfy $\dot u_i=H_i(u)-f_i(m)$, the forward equation, and $u(1)=g(m(1))$. Its potential objective is zero.

For a competing admissible control set $a_{21}=a_{31}=c=\log2$ and all other off-diagonal rates to zero. Then
\[
 m_2(t)=m_3(t)=\tfrac13e^{-ct},\qquad m_1(t)=1-\tfrac23e^{-ct}.
\]
At time one the law is $(2/3,1/6,1/6)$. Its energy is $c/6$ and its terminal potential is $-1/3$, giving
$J=(\log2)/6-1/3<0$.
For the last strict inequality, $\log2<1<2$, for example by integrating $1/x<1$ over $(1,2)$. Thus the stationary equilibrium path cannot minimize $J$.
\end{proof}

\paragraph{Computation.}
The competitor's objective evaluates to approximately $-0.2178088032$ in the script; the exact proof uses the analytic formula. For the admissible constant-repulsion example with $d=3$, $m_0=(0.5,0.3,0.2)$ and $T=1$, it evaluates the exact expectation formula at $\varepsilon=0.1,0.01,0.001$, obtaining first-coordinate means approximately $0.424802,0.490294,0.499003$. This is a deterministic formula check, not a simulated proof of tightness or of general selection.
\begin{lstlisting}
for eps in [0.1, 0.01, 0.001]:
    mean = [1/d + (x-1/d)*exp(-2*d*eps*T) for x in m0]
    assert abs(sum(mean)-1) < 1e-12
competitor_J = log(2)/6 - 1/3
assert competitor_J < 0
\end{lstlisting}

\subsection*{5) VERIFICATION}
The noise bracket contains the factor two generated by $dW^{ij}-dW^{ji}$; replacing it by independent coordinate Brownian noise would invalidate the proof. Localization uses only local repulsion convergence and a fixed interior initial law. The explicit family's PDE uniqueness and interior SDE existence were both checked, rather than treating the adjective ``admissible'' as a proof. The negative-potential competitor is finite energy and solves its forward equation, but it is \emph{not} claimed to be a noisy limit. Accordingly it refutes only ``every deterministic equilibrium is a global minimizer.'' It does not refute the input's selection claim. No corrected running costs have been introduced anywhere in the positive theorem.

\Needspace{8\baselineskip}
\subsection*{6) FINAL}
\textbf{LABEL: UNRESOLVED} for arbitrary smooth potentials and all admissible repulsion families.

\textbf{Strongest checked partial result.} For zero potentials, the full selected population path converges in probability to the unique minimizing path for every admissible family, with an explicit stopped-path error estimate. The constant family $2\varepsilon$ is independently proved admissible. An explicit smooth three-state example proves that identifying a limit as a deterministic equilibrium would not, by itself, finish the general problem.

\textbf{Exact first gap.} No inequality has been proved forcing the potential objective of a general uncorrected noisy limit to equal the global infimum of $J$; deterministic equilibrium equations alone do not imply it.

\textbf{Top three next moves.} (1) Derive an exact uncorrected It\^o energy identity and isolate every term lacking a potential sign. (2) Obtain uniform interior/tightness estimates for nonzero controlled drifts, stating the repulsion width and height explicitly. (3) Test the displayed concave terminal-potential example by solving its actual parabolic system and checking whether any admissible limit selects its nonminimizing constant equilibrium.

\textbf{Minimal counterexample perspective.} A three-state nonconvex potential with multiple deterministic equilibria is a natural test bed. A genuine counterexample must additionally supply an admissible repulsion family, prove convergence of the actual noisy law, and verify positive limiting mass on nonminimizers. Those additional claims have not been established here.

\paragraph{Reproducibility and scope.} This is one of eleven companion reports. The bundle includes \texttt{verify.py} and its executed results. Finite experiments supplement the proofs; they are not proofs of asymptotic claims. Literature coverage is targeted, and no novelty or formal proof-assistant certification is claimed.

\begin{thebibliography}{99}
\small
\bibitem{cardaliaguet-delarue}
Pierre Cardaliaguet and Fran\c{c}ois Delarue.
\emph{Selected topics in mean field games}. Proceedings of the International Congress of Mathematicians 2022, volume 5.
DOI: 10.4171/ICM2022/17. Sections 4.3--4.4; especially printed pp.~3691--3693.
\url{https://ems.press/content/book-chapter-files/33265}.

\bibitem{cecchin-delarue}
Alekos Cecchin and Fran\c{c}ois Delarue.
Selection by vanishing common noise for potential finite state mean field games.
\emph{Communications in Partial Differential Equations} 47(1) (2022), 89--168.
DOI: 10.1080/03605302.2021.1955256. Preprint: arXiv:2005.12153.
\url{https://arxiv.org/abs/2005.12153}.

\end{thebibliography}

\end{document}
